\documentclass[a4paper,12pt]{article}

\usepackage[margin=5mm]{geometry} % setup small page margins
\usepackage[czech]{babel} % add support for czech symbols
\usepackage[utf8]{inputenc} % input files are UTF-8 encoded
\usepackage[T1]{fontenc} % output font is encoded using T1
\usepackage{multicol} % allow multicols command for 2-column pages
\usepackage{hyperref} % make references clickable
\usepackage{listings} % lstlisting command to show tabs

% input enconding in UTF-8 may contain many symbols unknown to lstlisting
% what follows is a mapping of UTF-8 to LaTeX code
\lstset{literate=%
{á}{{\'a}}1
{é}{{\'e}}1
{í}{{\'i}}1
{ó}{{\'o}}1
{ú}{{\'u}}1
{ů}{{\r u}}1
{ý}{{\'y}}1
{ž}{{\v z}}1
{š}{{\v s}}1
{č}{{\v c}}1
{ř}{{\v r}}1
{ě}{{\v e}}1
{ď}{{\v d}}1
{ť}{{\v t}}1
{ň}{{\v n}}1
{Á}{{\'A}}1
{É}{{\'E}}1
{Í}{{\'I}}1
{Ó}{{\'O}}1
{Ú}{{\'U}}1
{Ů}{{\r U}}1
{Ý}{{\'Y}}1
{Ž}{{\v Z}}1
{Š}{{\v S}}1
{Č}{{\v C}}1
{Ř}{{\v R}}1
{Ě}{{\v E}}1
{Ď}{{\v D}}1
{Ť}{{\v T}}1
{Ň}{{\v N}}1
{â}{{\^a}}1
{’}{{'}}1
{-}{{-}}1
{“}{{`}}1
{”}{{'}}1
{´}{{'}}1
{ }{{ }}1
{—}{{--}}1
{–}{{--}}1
{...}{{\dots}}1
}

\lstset{%
    numberstyle=\phantom,%hidden
    includerangemarker=false,
    % inputencoding=UTF-8,
    basicstyle=\ttfamily\footnotesize,
    breakatwhitespace=true,
    keepspaces=true,
    showspaces=false,
    columns=fullflexible,
    frame=single,
    numbers=left,
    stepnumber=1,
    breaklines=true,
    tabsize=4
}

\begin{document}

\begin{multicols}{2}
    \subsection*{Kytarový zpěvník}
    \hrule\bigskip
    V ruce držíš \textit{první vydání (\today)} kompilace klasických a moderních písní s kytarovými akordy k ohni či do hospody.
    Písně jsou zaznamenány standardním způsobem 
    s těmito zásadními úpravami:
    \begin{itemize}
        \setlength\itemsep{-1mm}
        \item Text je psán ve větách, včetně čárek a teček.
        \item Mezery ve zpívání mají větší mezeru v textu.
        \item Důraz na slabice předchází pomlčka (\texttt{-}).
        \item Opakování stejného akordu je vždy vyznačeno opakovaným znakem.
        \item Hraní akordu celou dobu značíme jen názvem (\texttt{C}).
        \item Akord hraný kratší dobu u sebe má navíc znak \texttt{C,} je $\frac 34$ doby; \texttt{C.} je $\frac 12$ doby; \texttt{C'} je $\frac 14$ doby nebo pouze jednou zahraný akord.
    \end{itemize}

    Další informace a zdojové kódy lze najít na stránce \texttt{https://github.com/vaclavblazej/tabs}

    \subsection*{Jak byl zpěvník vytvořen}
    Po několik let byly střádány přejaté písně z internetu kam jsem zavedl opravy.
    Seznam písní je přejat z veřejného repozitáře, kde jsou písně uchovávány v surovém stavu.
    Potom jsou zpracovány programem a automaticky je vytvořen seznam podle názvů a interpretů.
    Zdrojové kódu všech programů, které se na tomto procesu podílí je zveřejněn na Github.
\end{multicols}
\newpage

%%% LIST %%%

%%% CONTENT %%%

\end{document}
